\section{Introduction}
\label{sec:intro}

Systematic reviews answer questions about the effects of health interventions by finding, appraising and combining
every study that meets criteria set in advance \citep{higgins2023handbook,page2021prisma}. Finding and selecting
those studies is the most labour-intensive part: searches across several databases return thousands of records,
and two reviewers screen each one \citep{borah2017analysis,lefebvre2023searching}. Automating parts of this work
has been studied for two decades, mostly as text classification that ranks or filters records for human screeners
\citep{cohen2006reducing,omara2015using,marshall2019toward,vandeschoot2021open}.

Large language models (LLMs) have widened the scope of such tools from ranking records to writing search strings,
screening titles and abstracts, extracting data and writing the synthesis \citep{wang2023can,guo2024automated,
khraisha2024can,asai2024openscholar}. They also bring a problem that ranking models did not: the same question can
produce a different set of studies each time it is asked, because the model's plan, the database responses and the
model's judgements all vary between runs. Repeating the same data extractions with a commercial LLM research
assistant agreed on 90\% of the extracted values but on only 46\% of the supporting quotations and 30\% of the
reasoning \citep{lagisz2026elicit}. A review whose
included studies cannot be reproduced cannot be audited or updated, and recent guidance on the responsible use of
AI in evidence synthesis asks tool developers to report how their tools were evaluated, to keep people in control
of decisions, and to label tools without robust evaluation as such \citep{raise2026,jointstatement2025}.

Most evaluations of LLM-based tools report agreement with human screening decisions on a fixed set of records
\citep{guo2024automated,khraisha2024can}. Fewer measure the whole path from question to included studies: how many
of the studies a published review included would the tool have found, at which step are the others lost, and does
a repeat of the same run give the same result? These are the questions a reviewer needs answered before relying on
such a tool for a search, and they have a clear reference standard in the included studies of published
systematic reviews \citep{kanoulas2019clef}.

This paper describes \tool{}, an open-source tool that separates the two kinds of work. The parts where a model
helps but may vary (planning searches, screening against criteria, writing) are bounded and recorded, and the part
that decides which studies are included (merging duplicates, eligibility and ranking) follows fixed, published
rules applied to stored records, so that it can be replayed exactly. Our contributions are:
\begin{enumerate}
\item A \emph{protocol-first, replayable} design for LLM-assisted study selection: the review protocol is set
  before searching; search strings are compiled by fixed rules; the rule-based selection does not read the model's
  plan; every record, decision and model judgement is stored in a corpus with a content hash, from which the
  selection can be reproduced without searching again or calling a model (Section~\ref{sec:design}).
\item A \emph{pre-registered evaluation} on 33 Cochrane intervention reviews with a development and held-out test
  split, in which the measured code was frozen before test results were used, reporting recall at each stage,
  where included studies are lost, baselines, and the repeatability of each step (Sections~\ref{sec:eval}
  and~\ref{sec:results}).
\item An analysis of the trade-off this design makes: keeping the selection independent of the model makes it
  reproducible but costs recall when the rule-based reading of the question is poor
  (Section~\ref{sec:discussion}).
\end{enumerate}
